\chapter{Basic Analysis}

\section{Real Numbers}
\begin{definition}
    Let $\mathcal{R} \defequal \{ \{a_n\} \mid \text{$\{a_n\}$ is Sequence of Rational numbers} \}$. \\
    Define an Equivalent Relation on $\mathcal{R}$ as: for all $\{a_n\}, \{b_n\} \in \mathcal{R}$,
    \begin{center}
        $\{a_n\} \sim \{b_n\}$ \textit{if and only if} for any $\epsilon \in \mathbb{Q}^+$, there exists $N \in \mathbb{N}$ such that $n \geq N \implies |a_n - b_n| < \epsilon$.
    \end{center}
\end{definition}

\section{Basic Topology on $\mathbb{R}$}
\subsection{Representation Theorem for open sets on $\mathbb{R}$}
\begin{definition}
    Let $S$ be an open set of $\mathbb{R}$.\\
    An open interval $I$ is called a \textit{component interval} of $S$ if: $I \subseteq S$ and
    \begin{center}
        There exists no open interval $J$ such that $I \subsetneq J \subseteq S$
    \end{center}
    Note: Component interval is not unique; $S = (0,1) \cup (1,2)$ has two component intervals.
\end{definition}
\begin{lemma}
    
\end{lemma}

\section{Sequence}

\section{Series}

\section{Differentiation}

\section{Riemann Integral}

\section{Sequence of Functions}
